% Master technical CV (engineering track). Layout: templates/cv.cls.
% Replace the example content with your own, or let Claude build it from your profile
% ("build my base CVs from the profile"). Applications copy this file and tailor it.
% End every bullet with "% src: <entry-id>#F<n>" so each claim traces to a profile fact.
\documentclass{cv}

\cvname{{{NAME}}}
\cvlocation{{{CITY}}}
\cvemail{{{EMAIL}}}
\cvphone{{{PHONE}}}{{{PHONE_LINK}}}
\cvlinkedin{{{LINKEDIN_URL}}}{{{NAME}}}

\begin{document}

\makecvheader

\section{Key Highlights}
\begin{bullets}
  \item M.Sc. student in <field> at <university> with hands-on experience from <job or organisation>. % src: example#F1
  \item <Your strongest technical claim, one line.> % src: example#F2
\end{bullets}

\section{Education}
\education{<University>}{<Degree> (<year> of <n> years)}{Aug 2024 -- Present}
\begin{bullets}
  \item \textbf{Average grade:} <grade> % src: example#F3
  \item \textbf{Coursework:} <three to five relevant courses>. % src: example#F4
\end{bullets}
\educationshort{<Upper secondary school>}{<Programme>}{Aug 2021 -- June 2024}

\section{Technical Experience}
\entry{<Job title>}{<Organisation>}{<City, Country>}{June 2026 -- Aug 2026}
\begin{bullets}
  \item <What you built or did, with a concrete detail.> % src: example#F5
  \item <The result, or what it was used for.> % src: example#F6
\end{bullets}

\section{Projects}
\begin{bullets}
  \item \textbf{<Project>:} <one line on what it does and your part>. % src: example#F7
\end{bullets}

\section{Skills}
\begin{skills}
  \skill{Programming}{Python, C++, MATLAB}
  \skill{Tools}{Linux, Git, Docker}
  \skill{Languages}{<language> (native), English (fluent)}
\end{skills}

\section{Awards}
\begin{bullets}
  \item <Competition, placement, year.> % src: example#F8
\end{bullets}

\end{document}
